% fduthesis: 复旦大学博士学位论文
% 使用 XeLaTeX 编译

\documentclass[type=doctor]{fduthesis}

\fdusetup{
  style = {
    footnote-style = xits,
    bib-backend = bibtex,
    bib-resource = {main.bib},
  },
  info = {
    title = {数据选择驱动的深度学习效率与可靠性优化研究},
    title* = {Data Selection Driven Optimization of \\
              Deep Learning Efficiency and Reliability},
    author = {万玮霖},  % TODO: 确认姓名
    supervisor = {TODO \quad 教授},  % TODO: 填写导师信息
    major = {计算机科学与技术},  % TODO: 确认专业
    degree = academic,
    department = {计算机科学技术学院},  % TODO: 确认院系
    student-id = {TODO},  % TODO: 填写学号
    keywords = {数据选择, 深度学习, 核心子集选择, 模型剪枝, 分布外检测},
    keywords* = {Data Selection, Deep Learning, Coreset Selection, Model Pruning, Out-of-Distribution Detection},
    clc = {TP181},  % TODO: 确认分类号
  }
}

% ======== 宏包 ========
\usepackage{amsmath, amssymb, amsfonts}
\usepackage{bm}           % 粗体数学符号
\usepackage{booktabs}     % 三线表
\usepackage{multirow}     % 合并行
\usepackage{algorithm}
\usepackage{algorithmic}
\usepackage{graphicx}
\usepackage{subcaption}   % 子图

% ======== 自定义命令 ========
% 通用
\newcommand{\btheta}{\bm{\theta}}
\newcommand{\bm}[1]{\boldsymbol{#1}}

% CADS 专用
\newcommand{\Lval}{\mathcal{L}_{\text{val}}}
\newcommand{\Ltrn}{\mathcal{L}_{\text{trn}}}
\newcommand{\thetaC}{\btheta_C(\mathbf{m})}

% SWaST 专用
\newcommand{\Lce}{\mathcal{L}_{\text{CE}}}
\newcommand{\Lsp}{\mathcal{L}_{\text{SP}}}
\newcommand{\Ltotal}{\mathcal{L}_{\text{total}}}

% NAP 专用
\newcommand{\SNAP}{S_{\text{NAP}}}
\newcommand{\SNAPformer}{S_{\text{NAP}}^{\text{Former}}}

\begin{document}

% ======== 前置部分 ========
\frontmatter

\tableofcontents

\begin{abstract}
深度学习的突破性进展高度依赖大规模高质量数据，但数据规模增长带来的计算成本攀升、标注资源稀缺和数据质量参差等问题日益严峻。数据选择——在深度学习的不同阶段，根据特定目标从候选数据中识别和利用最相关样本的过程——是应对这些挑战的核心方法论。然而，现有研究大多将数据选择作为孤立的优化问题，忽视了选择策略与计算资源、模型容量、推理时分布偏移等关键因素之间的深层交互。

本文围绕"深度学习全流程数据选择"这一主线，从阶段视角（训练$\to$推理）和交互变量视角（计算资源$\to$模型容量$\to$分布偏移）两个互补维度出发，系统研究数据选择与上述关键因素的交互，提出三项工作：

针对训练阶段数据选择与计算资源约束的交互，提出计算预算感知的数据选择框架CADS。CADS首次将计算预算作为双层优化中的显式约束，通过概率化重参数化的策略梯度绕过内层非收敛的困难，并通过惩罚松弛与一维损失近似将双层优化转化为高效的单层求解。CADS揭示了一个被广泛忽视的现象：最优数据子集是计算预算的函数——低预算偏好简单干净的数据，高预算受益于多样化的数据。在视觉分类和语言模型微调任务上，CADS相较基线方法最高提升14.42\%。该工作发表于NeurIPS 2025。% [VERIFY]

针对训练阶段数据选择与模型容量约束的交互，提出数据选择与权重剪枝的协同优化框架SWaST。SWaST首次透明地分析了数据冗余与参数冗余的双向耦合关系——冗余样本阻碍有效剪枝，冗余参数干扰数据选择——并通过交替执行权重剪枝与核心子集选择建立协同效应。SWaST识别并解决了同步优化中独有的"双重损失"问题，通过状态保护机制防止不可逆退化。在四个基准数据集和两种网络架构上，SWaST实现了最高17.83\%的准确率提升，同时削减10\%至90\%的浮点运算量。该工作发表于AAAI 2026。% [VERIFY]

针对推理阶段数据选择与分布偏移的交互，提出基于神经激活先验的分布外检测方法NAP。NAP发现了一个被现有方法忽视的检测信号：在全局池化前的特征图中，分布内样本在各通道内呈现尖锐的空间激活模式，而分布外样本的激活更为分散。基于这一先验，NAP设计了即插即用的评分函数，同时适用于CNN和Vision Transformer架构，并可与现有方法通过加权几何均值实现正交互补。在CIFAR和ImageNet基准上，NAP与多种基线方法组合后FPR95最高降低66.03\%。% [VERIFY]

三项工作从训练侧延伸到推理侧，共同揭示了一个方法论启示：数据选择的最优策略并非孤立决定的，而是取决于其所处的约束环境——计算资源约束改变最优训练数据的组成，模型容量约束与数据冗余需要协同而非独立处理，推理时分布偏移要求对数据有效性做在线判断。
\end{abstract}

\begin{abstract*}
The breakthroughs in deep learning are highly dependent on large-scale, high-quality data, yet the growth of data scale brings increasingly severe challenges including escalating computational costs, scarcity of annotation resources, and inconsistent data quality. Data selection---the process of identifying and utilizing the most relevant samples from candidate data at different stages of deep learning---is a core methodology for addressing these challenges. However, existing research largely treats data selection as an isolated optimization problem, overlooking the deep interactions between selection strategies and key factors such as computational resources, model capacity, and inference-time distribution shift.

This thesis centers on ``data selection throughout the deep learning pipeline'' and systematically studies the interactions between data selection and the aforementioned key factors from two complementary perspectives: the stage perspective (training $\to$ inference) and the interaction variable perspective (computational resources $\to$ model capacity $\to$ distribution shift). Three contributions are presented:

For the interaction between data selection and computational resource constraints during training, we propose CADS (Computation-Aware Data Selection), a framework that, for the first time, incorporates the computational budget as an explicit constraint in bilevel optimization for data selection. CADS bypasses the difficulty of inner-level non-convergence through probabilistic reparameterization with REINFORCE policy gradients, and further transforms the bilevel optimization into an efficient single-level problem via penalty relaxation and one-dimensional loss approximation. CADS reveals a widely overlooked phenomenon: the optimal data subset is a function of the computational budget---low budgets favor simple, clean data while high budgets benefit from diverse data. On visual classification and language model fine-tuning tasks, CADS achieves up to 14.42\% improvement over baseline methods. This work was published at NeurIPS 2025.

For the interaction between data selection and model capacity constraints during training, we propose SWaST (Simultaneous Weight and Sample Tailoring), a framework for synergistic optimization of data selection and weight pruning. SWaST is the first to transparently analyze the bidirectional coupling between data redundancy and parameter redundancy---redundant samples hinder effective pruning, while redundant parameters impair data selection---and establishes synergy through alternating weight pruning and coreset selection. SWaST identifies and addresses the unique ``Critical Double-Loss'' problem in simultaneous optimization through a state protection mechanism that prevents irreversible degradation. On four benchmark datasets with two network architectures, SWaST achieves up to 17.83\% accuracy improvement while reducing 10\% to 90\% of FLOPs. This work was published at AAAI 2026.

For the interaction between data selection and distribution shift during inference, we propose NAP (Neural Activation Prior), an out-of-distribution (OOD) detection method based on a previously overlooked detection signal: in the feature maps before global average pooling, in-distribution samples exhibit sharp spatial activation patterns within each channel, while OOD samples show more dispersed activations. Based on this prior, NAP designs a plug-and-play scoring function applicable to both CNNs and Vision Transformers, achieving orthogonal complementarity with existing methods through weighted geometric mean combination. On CIFAR and ImageNet benchmarks, NAP reduces FPR95 by up to 66.03\% when combined with various baseline methods.

Together, the three works span from training to inference and reveal a methodological insight: the optimal data selection strategy is not determined in isolation, but depends on its operating constraints---computational resource constraints alter the composition of optimal training data, model capacity constraints and data redundancy require synergistic rather than independent treatment, and inference-time distribution shift demands online assessment of data validity.
\end{abstract*}

% ���号��
\begin{notation}[lp{10cm}]
  $\btheta$                   & 模型参数 \\
  $\mathcal{D}$               & 完整训练数据集 \\
  $\mathcal{L}$               & 损失函数 \\
  $f$                         & 神经网络模型 \\
  $x, y$                      & 输入样本、标签 \\
  $\mathbf{m}$                & 二值数据选择向量（样本掩码） \\
  $\mathbf{s}$                & 可学习的采样分布参数 \\
  $C$                         & 计算预算（第三章）/ 通道数以$N_c$表示（第五章） \\
  $\hat{\mathcal{D}}$         & 核心子集（coreset） \\
  $\mathcal{R}$               & 核心子集选择的epoch间隔 \\
  $\lambda$                   & 状态保护损失权重 \\
  $\mathbf{A}$                & 全局池化层前的激活张量 \\
  $S_{\text{NAP}}$            & NAP评分函数 \\
\end{notation}

% ======== 主体部分 ========
\mainmatter

\chapter{绪论}

\section{研究背景：深度学习对数据的依赖与数据选择的核心地位}

过去十余年间，深度学习在计算机视觉、自然语言处理和多模态理解等领域取得了一系列里程碑式的突破。2012年，AlexNet在ImageNet大规模视觉识别挑战赛\cite{deng2009imagenet}中的优异表现标志着深度学习时代的开启，此后卷积神经网络（Convolutional Neural Network, CNN）迅速成为图像分类、目标检测和语义分割等视觉任务的主流方法\cite{he2016deep, simonyan2015very}。在自然语言处理领域，以GPT系列\cite{radford2019language}和Transformer\cite{vaswani2017attention}为代表的大规模预训练语言模型（Large Language Model, LLM）证明了通过海量文本数据训练的模型可以涌现出强大的语言理解与生成能力。在生成模型方面，扩散模型（Diffusion Model）在图像生成任务上展现了前所未有的质量与多样性。这些突破虽然涉及不同的模型架构和训练范式，但都建立在一个共同的基础之上——大规模高质量数据。

在大模型时代，数据的地位进一步凸显。研究者们逐渐认识到，LLM的许多能力涌现并非仅源于模型规模的增长或架构的改进，而是在很大程度上取决于训练数据的规模、质量与多样性\cite{hoffmann2022training, kaplan2020scaling}。高质量数据的获取与筛选已成为前沿AI研究与工业实践中的核心竞争力\cite{albalak2024survey}。Ilya Sutskever曾发出"数据即将耗尽"的警示——无论这一判断��终是否准确，它深刻反映了当前研究界对数据资源的高度依赖与焦虑。这一趋势也暗示着，单纯追求"更多数据"的粗放式路径正在接近其极限，"更聪明地使用数据"正变得愈发迫切。

然而，数据规模的增长在带来模型能力提升的同时，也引发了一系列现实挑战。首先，计算成本随数据量的增长呈超线性增长\cite{hoffmann2022training}——训练一个大规模模型所需的计算资源已达到普通研究机构难以承受的水平。其次，高质量标注数据日益稀缺且获取成本高昂，尤其在医疗、法律等专业领域，专家标注的成本使得全量标注变得不切实际。此外，大量收集的数据中不可避免地存在冗余、噪声和分布不均等问题，这些低质量数据不仅浪费计算资源，还可能损害模型的泛化性能\cite{sorscher2022beyond}。上述挑战共同指向一个核心认识：\textbf{并非所有数据都同等有价值}。如何从大量候选数据中识别和利用最相关的数据，成为提升深度学习效率与可靠性的关键问题。

数据选择（Data Selection）正是应对这一挑战的核心方法论\cite{albalak2024survey}。广义而言，数据选择是指在深度学习的不同阶段，根据特定目标从候选数据中识别、筛选或评估最相关样本的过程。它不仅涵盖训练阶段常见的核心子集选择（Coreset Selection）\cite{mirzasoleiman2020coresets, killamsetty2021grad}——即从全量训练数据中选取具有代表性的子集用于高效训练，还包括数据混合配比（Data Mixture）——即确定不同来源数据的最优混合比例，数据质量过滤——即剔除低质量或有害样本\cite{chen2023alpagasus, zhou2023lima}，以及推理阶段对输入数据有效性的判断——即判定模型是否应对某一输入做出预测。值得指出的是，数据选择与数据增强（Data Augmentation）���数据集蒸馏（Dataset Distillation）\cite{wang2018dataset, cazenavette2022dataset}有着本质区别：数据选择是从现有数据中"选"，而后两者是通过变换或优化"造"新数据。本文聚焦于数据选择这一方向。

数据选择并非仅局限于深度学习流程中的某一特定环节，而是以不同形式贯穿训练、微调与推理等多个阶段。在\textbf{训练阶段}，核心任务是从全量数据中选出最具信息量的子集，在有限计算资源下最大化模型性能\cite{mirzasoleiman2020coresets, paul2021deep}。在\textbf{微调阶段}，面对来源多样、质量参差的候选数据（如人工标注数据、模型生成数据、开源数据集等），需要选择最匹配下游任务的数据子集或配比，在有限微调预算下高效适配预训练模型\cite{zhou2023lima, chen2023alpagasus}。在\textbf{推理阶段}，模型部署到真实环境后，输入数据的来源不再可控，需要判断每一个输入样本是否处于模型的"有效数据域"内——即模型是否有能力对该输入做出可靠预测\cite{hendrycks2016msp, yang2021generalized_survey}。数据选择在这三个阶段以不同的形式出现、面临不同的约束条件，但都服务于同一个目标：让深度学习系统更高效、更可靠地运行。

进一步地，数据选择的最优策略并非孤立决定的，而是与多种关键因素存在深层交互。\textbf{计算资源约束}影响最优训练数据子集的组成——在不同计算预算下，对模型最有价值的数据可能截然不同\cite{yin2024compute}。\textbf{模型容量约束}与数据冗余之间存在耦合——当模型同时进行压缩（如剪枝\cite{han2015learning}）时，数据的最优选择策略需要随模型状态的变化而调整。\textbf{推理时分布偏移}则要求在推理阶段对输入数据的有效性做出在线判断——训练时的数据选择界定了模型的有效域边界，推理时需要检测输入是否越出这一边界\cite{yang2021generalized_survey}。然而，现有的数据选择研究大多将问题孤立处理：训练数据选择方法不考虑计算预算的影响\cite{killamsetty2021grad, paul2021deep}，数据选择与模型压缩被视为独立的流水线步骤\cite{evci2020rigging}，推理时的分布外检测与训练时的数据选择之间缺乏联系\cite{sun2021react, liu2020energy}。本文旨在填补这一空白，系统研究数据选择与上述关键因素的交互。


\section{研究问题}

基于上述分析，本文从两个互补的视角出发，围绕数据选择提出三个核心研究问题。\textbf{阶段视角}（主线A）关注数据选择在深度学习不同阶段扮演的角色，\textbf{交互变量视角}（主线B）关注数据选择与关键因素之间的交互如何影响选择策略。表~\ref{tab:research-questions}展示了三个研究问题在这两个视角下的定位。

% [表1-1: 三个研究问题的双视角定位]
\begin{table}[htb]
  \centering
  \caption{三个研究问题的双视角定位}
  \label{tab:research-questions}
  \begin{tabular}{cll}
    \toprule
    研究问题 & 主线A：阶段视角 & 主线B：交互变量视角 \\
    \midrule
    子问题一 & 训练/微调阶段的数据选择 & 数据选择 $\times$ 计算资源约束 \\
    子问题二 & 训练阶段的数据与模型联合选择 & 数据选择 $\times$ 模型容量约束 \\
    子问题三 & 推理阶段的数据相关性判断 & 数据选择 $\times$ 推理时分布偏移 \\
    \bottomrule
  \end{tabular}
\end{table}

三个子问题呈递进关系：从阶段视角看，从训练侧（子问题一、二）延伸到推理侧（子问题三）；从交互视角看，交互维度从计算资源到模型容量再到推理时分布偏移，逐步扩展。

\subsection{子问题一：计算预算约束下的训练数据选择}

第一个研究问题聚焦于训练与微调阶段最基础的数据选择问题：\textbf{给定有限的计算预算，如何从全量���练数据中选出使模型性能最优的子集？}

从阶段视角看，这是数据选择在训练阶段的直接体现。从交互视角看，这一问题的核心在于数据选择与计算资源约束之间的交互。已有研究表明，不同计算预算下最优数据子集的组成截然不同\cite{sorscher2022beyond, yin2024compute}：当计算预算有限时，模型仅能学习数据中的简单模式，此时简单、干净的样本更有价值；当计算预算充裕时，模型有能力从多样化甚至含噪声的数据中提取更丰富的信息。然而，现有的数据选择方法普遍忽视计算预算这一关键变量——它们隐含地假设最优数据子集独立于训练时间\cite{killamsetty2021grad, paul2021deep, toneva2018empirical}，导致没有任何单一算法能在不同计算预算下一致领先\cite{xia2024rethinking}。这一问题的核心挑战在��：如何将计算预算作为数据选择框架中的显式变量，使选择策略能够自适应地随预算变化。

\subsection{子问题二：数据选择与模型剪枝的协同}

第二个研究问题同属训练阶段，但将"选择"的维度从数据扩展到模型：\textbf{当训练数据选择与模型剪枝同步进行时，如何利用二者的交互实现协同优化？}

从阶段视角看，模型剪枝（Weight Pruning）是训练阶段的另一种效率优化手段，它通过移除冗余参数来压缩模型\cite{han2015learning, lecun1989optimal}。从交互视角看，数据冗余与参数冗余之间存在双向耦合：一方面，冗余样本（如噪声数据）迫使模型为拟合这些样本而保留不必要的权重，从而阻碍有效剪枝；另一方面，冗余参数扩大了模型的有效表示空间，增加了核心子集选择的搜索难度，导致选出的子集质量下降。然而，现有工作将数据选择和模型剪枝视为独立的流水线步骤——先选数据再训练完整模型，或先训练完整模型再剪枝——未利用也未分析二者之间的耦合关系\cite{evci2020rigging, killamsetty2021grad}。这一问题的核心挑战在于：如何在训练过程中同时进行数据选择与模型剪枝，利用它们的交互建立协同效应。

\subsection{子问题三：推理时的数据有效性判断}

第三个研究问题将数据选择的视角从训练阶段延伸至推理阶段：\textbf{模型部署后，如何判断输入样本是否处于训练数据选择所界定的"有效数据域"内？}

从阶段视角看，模型训练完成并部署后，输入数据的来源不再可控。当输入样本偏离训练数据分布时，模型可能产生过度自信的错误预测\cite{hendrycks2016msp}，这在安全关键应用（如自动驾驶\cite{filos2020can, janai2020computer}、医学诊断\cite{pooch2020can}）中尤为危险。分布外检测（Out-of-Distribution Detection, OOD Detection）可被理解为推理时的数据"选择"——选择（接受）属于训练分布内的输入，拒绝偏离训练分布的输入\cite{yang2021generalized_survey}。从交互视角看，训练阶段的数据选择界定了模型的有效域边界；当推理时输入分布发生偏移，需要判断输入是否仍属于该有效域。这是数据选择与推理时分布偏移（Inference-time Distribution Shift）的交互。这一问题的核心挑战在于：现有的后处理OOD检测方法主要利用全局平均池化后的特征或最终输出层信息\cite{liu2020energy, hendrycks2016msp, sun2021react, sun2022dice}，丢失了池化前通道内的空间激活分布信息，限制了检测能力。


\section{本文贡献概览}

针对上述三个研究问题，本文围绕数据选择这一核心方法论，分别提出一个框架或方法，共同构成对数据选择问题从训练到推理的系统性研究。本文的主要贡献如下：

\textbf{贡献一：计算预算感知的数据选择框架CADS（第三章）。}针对子问题一，本文提出CADS（Computation-Aware Data Selection）框架，首次将计算预算作为双层优化中的显式约束纳入数据选择\cite{zhou2022probabilistic}。在技术层面，CADS通过概率化重参数化和REINFORCE策略梯度实现无Hessian的外层梯度估计，并进一步提出基于惩罚松弛的单层目标转化，将计算开销降低至标准双层优化的3至20倍以下。% [VERIFY]
CADS包含两种变体：CADS-E实现样本级选择，CADS-S实现来源级选择，在视觉分类和语言模型微调任务上，相较基线方法最高提升14.42\%。% [VERIFY]
该工作发表于NeurIPS 2025。

\textbf{贡献二：数据选择与权重剪枝的协同优化框架SWaST（第四章）。}针对子问题二，本文提出SWaST（Simultaneous Weight and Sample Tailoring）框架，通过交替执行权重剪枝与核心子集选择建立协同效应。SWaST首次透明地分析了数据冗余与参数冗余的双向耦合关系，发现并定义了"双重损���"（Critical Double-Loss）问题——即关键参数与其支撑样本被同时移除导致的不可逆退化。为解决这一问题，SWaST引入了状态保护机制，通过知识蒸馏约束防止协同优化过程中的灾难性遗忘。在四个基准数据集和两种网络架构上，SWaST实现了最高17.83\%的准确率提升，同时削减10\%至90\%的浮点运算量（FLOPs）。% [VERIFY]
该工作发表于AAAI 2026。

\textbf{贡献三：基于神经激活先验的分布外检测方法NAP（第五章）。}针对子问题三，本文提出NAP（Neural Activation Prior），发现了一个被现有方法忽视的OOD检测信号：在全连接层前的特征图中，分布内（In-Distribution, ID）样本在各通道内倾向于激活少量神经元且产生较大激活值，而分布外样本的激活则更为分散。基于这一神经激活先验，NAP设计了即插即用的评分函数，无需额外训练，并可与现有OOD检测方法\cite{liu2020energy, hendrycks2016msp, sun2021react, sun2022dice, djurisic2022extremely, sun2022knn}通过加权几何均值进行正交互补。在CIFAR\cite{krizhevsky2009cifar}和ImageNet\cite{deng2009imagenet}基准上，NAP与多种基线方法组合后，FPR95最高降低66.03\%。% [VERIFY]
NAP同时支持CNN和Vision Transformer\cite{dosovitskiy2020image}架构。该工作目前处于投稿中。


\section{论文结构}

图~\ref{fig:research-map}展示了本文三项工作在双主线框架下的位置关系。横轴为阶段递进，从训练侧（第三、四章）延伸到推理侧（第五章）；纵轴为交互维度，从计算资源约束（第三章）到模型容量约束（第四章）再到推理时分布偏移（第五章）。主线A（阶段视角）回答"在深度学习的哪个阶段做数据选择"，主线B（交互变量视角）回答"数据选择与什么关键因素交互"。两条主线相互补充：主线A提供直觉清晰的叙事脉络，主线B揭示工作之间的深层统一性。

% [图1-1: 研究地图——三项工作在双主线框架下的位置。需用户手动绘制。]
\begin{figure}[htb]
  \centering
  % \includegraphics[width=0.8\textwidth]{figures/research-map.pdf}
  \fbox{\parbox{0.8\textwidth}{\centering\vspace{3cm}
    \textbf{占位图：研究地图}\\[0.5em]
    横轴：阶段递进（训练$\to$推理）\\
    纵轴：交互维度（计算资源$\to$模型容量$\to$分布偏移）\\
    三个方框分别对应CADS / SWaST / NAP
  \vspace{3cm}}}
  \caption{研究地图——三项工作在双主线框架下的位置关系}
  \label{fig:research-map}
\end{figure}

本文其余部分的组织如下：

\textbf{第二章\quad 相关工作}，从四个方面综述与本文研究相关的已有工作——训练数据选择方法、模型冗余与效率优化、数据选择中的优化方法、分布外检测，揭示现有研究忽视数据选择与关键因素交互的共性不足。

\textbf{第三章\quad 计算预算感知的训练数据选择}，研究数据选择与计算资源约束的交互。本章从一个探索性实验出发，揭示不同计算预算下最优数据子集的组成截然不同，进而提出将计算预算作为显式变量的数据选择框架CADS。

\textbf{第四章\quad 数据选择与权重剪枝的协同优化}，研究数据选择与模型容量约束的交互。本章首先通过多项式插值实验透明地分析数据冗余与参数冗余的双向耦合，然后提出交替优化框架SWaST，同时削减两类冗余并建立协同效应。

\textbf{第五章\quad 推理阶段的数据相关性判断}，研究数据选择与推理时分布偏移的交互。本章提出神经激活先验（NAP），利用全局池化前通道内的激活分布模式区分分布内与分布外样本，为推理时的数据有效性判断提供新的检测信号。

\textbf{第六章\quad 总结与展望}，从阶段覆盖和交互维度两个视角统一回顾三项工作的贡献，讨论各自的局限性，并展望数据选择领域的未来研究方向。

\chapter{相关工作}

本章围绕数据选择这一核心主题，从四个方面综述与本文研究密切相关的��有工作。第~\ref{sec:training-data-selection}节综述训练数据选择方法，涵盖核心子集选择、大模型微调中的数据选择以及相关概念（数据集蒸馏与课程学习），为第三、四章提供数据侧背景。第~\ref{sec:model-redundancy}节综述模型冗余与效率优化方法，聚焦结构化与非结构化剪枝以及稀疏训练，为第四章提供模型侧背景。第~\ref{sec:optimization-methods}节综述数据选择中常用的双层优化方法论，为第三章提供方法论工具。第~\ref{sec:ood-detection}节综述分布外检测方法，为第五章提供推理侧背景。本章的组织并非对四个独立领域的简单罗列，而是始终围绕"数据选择与关键因素的交互"这一视角进行串联。

需要指出的是，上述四个领域各自已有丰富的研究积累，但一个共性问题贯穿其中：现有工作倾向于将数据选择作为独立决策进行研究，较少关注数据选择与计算预算、模型容量、推理时分布偏移等因素之间的交互影响。本章在综述各领域时，会在每节末尾明确指出这一不足，为后续三章的研究动机提供铺垫。


\section{训练数据选择方法}\label{sec:training-data-selection}

\subsection{核心子集选择}

核心子集选择（Coreset Selection）旨在从完整训练集中识别出具有代表性的子集，使得在该子集上训练的模型能够尽可能逼近在全量数据上训练的性能\cite{albalak2024survey}。现有方法大致可分为三类：基于优化的方法、基于训练动态的启发式方法和基于几何距离的方法。

\textbf{基于优化的方法}通过显式优化目标来选择子集。子模函数最大化方法利用子模函数（如Facility Location函数）的近似最优性保证，在多项式时间内选出高质量子集\cite{wei2015submodularity, mirzasoleiman2016fast}。梯度匹配方法（如GradMatch\cite{killamsetty2021grad}）通过最小化子集梯度与全集梯度之间的差异来选择样本，使得在子集上的梯度更新尽可能模拟全集训练。双层优化方法（如PBCS\cite{zhou2022probabilistic}、GLISTER\cite{killamsetty2021glister}）将数据选择建模为外层问题，以验证集性能为目标优化子集组成，内层在所选子集上训练模型。这类方法通常具有较好的理论性质，但计算开销较高。

\textbf{基于训练动态的启发式方法}利用模型训练过程中产生的信号来估计样本重要性，无需额外的优化过程。Toneva等人\cite{toneva2018empirical}提出记录训练过程中每个样本被"遗忘"的次数（Forgetting Events），发现被反复遗忘的样本对模型学习至关重要。Paul等人\cite{paul2021deep}提出EL2N评分，利用训练早期的误差范数评估样本难度，以及GraNd评分利用梯度范数进行类似估计。Pruthi等人\cite{pruthi2020estimating}通过追踪梯度下降轨迹来估计训练数据的影响。这些方法计算高效，但缺乏选择质量的理论保障。

\textbf{基于几何距离的方法}在特征空间中根据样本间的距离关系选择代表性样本。K-Center-Greedy\cite{sener2018active}选择使得任何未选样本到最近已选样本的最大距离最小的子集。Moderate Coreset\cite{xia2022moderate}选择距离类中心适中的样本，避免过于简单或过于困难的极端样本。近年来，基于数据模型（Datamodels）的方法\cite{ilyas2022datamodels, engstrom2024dsdm}和基于最优传输的方法\cite{xiao2024feature}进一步丰富了这一领域。

\subsection{大模型微调中的数据选择}

随着大规模语言模型的兴起，微调阶段的数据选择面临新的挑战\cite{albalak2024survey}。与传统核心子集选择相比，LLM微调场景具有三个显著特殊性：数据来源多样（人工标注、模型生成、开源数据集），质量参差不齐，且单次训练成本极高\cite{rae2021scaling}，使得数据选择的经济价值尤为突出。

在这一背景下，一系列高效的数据选择方法应运而生。Zhou等人\cite{zhou2023lima}通过LIMA实验证明，仅需1000条精心筛选的高质量指令数据即可有效微调LLM，挑战了"数据越多越好"的传统认知。在自动化筛选方面，Chen等人\cite{chen2023alpagasus}提出利用模型自身对数据质量进行评估和过滤；Cao等人\cite{cao2023instruction}提出基于自然语言指标的数据挖掘方法；Lu等人\cite{lu2023instag}通过语义意图标签实现数据分类与筛选。Liu等人\cite{liu2023what}提出DEITA框架，综合数据质量、复杂度和多样性三个维度进行统一评估。近期工作还探索了基于梯度的技术路线，包括聚类核心子集选择\cite{zhang2024tagcos}和基于图的方法\cite{zhao2025beyond}，以及利用小模型训练轨迹指导大模型数据选择\cite{yang2024smalltolarge}和基于稀疏自编码器的多样性驱动选择\cite{yang2025diversity}等新思路。

与传统核心子集选择相比，LLM数据选择更依赖启发式质量指标而非梯度信息，且常需处理来源级（source-level）的数据混合配比问题——即确定不同来源数据的最优混合比例，而非逐样本选择。

\subsection{数据集蒸馏与课程学习}

除核心子集选择外，数据效率领域还有两个相关但不同的研究方向：数据集蒸馏和课程学习。

数据集蒸馏（Dataset Distillation）\cite{wang2018dataset}通过优化生成少量合成样本来压缩数据集的信息内容。代表性方法包括基于梯度匹配的方法\cite{cazenavette2022dataset}和基于可寻址记忆的方法\cite{deng2022remember}。与核心子集选择的"选"相对，蒸馏是"造"新数据\cite{yu2023dataset}。蒸馏方法的优点在于可实现极致的数据压缩率，但合成样本缺乏可解释性，且生成过程本身计算昂贵，在异构大规模数据上的扩展性有限。

课程学习（Curriculum Learning）\cite{bengio2009curriculum}按照从易到难的顺序组织训练样本，关注的是训练顺序而非子集选择\cite{hacohen2019power, soviany2021curriculum}。与核心子集选择有交叉之处——二者都涉及样本重要性评估——但目标不同：课程学习优化的是"以什么顺序使用所有数据"，而核心子集选择优化的是"使用哪些数据"。相关地，重要性采样\cite{alain2015variance, katharopoulos2018not, loshchilov2015online}和自适应重加权方法通过优先处理对模型改进贡献最大的数据点来提升训练效率。

本文聚焦于从现有数据中选择子集的问题，数据集蒸馏和课程学习不是本文的直接研究对象，但它们为理解数据效率的全景提供了有价值的参照。

\subsection{小结：忽视计算预算的交互}

上述所有训练数据选择方法在选择数据时，隐含假设训练会运行到收敛或至少运行足够长的时间，即最优子集与实际可用的计算预算无关。然而，实际应用中的计算预算差异巨大——从小规模实验的数千步到大规模预训练的数百万步。近期研究\cite{diddee2024chasing, shen2024rethinking, xia2024rethinking}揭示了一个令人意外的现象：精心设计的数据选择策略在某些计算预算下甚至不如随机选择。Yin等人\cite{yin2024compute}进一步证明，当计算预算被显式纳入考量时，此前有效的数据选择方法往往不再最优。这些发现暗示最优数据子集确实依赖于计算预算，而非一个与预算无关的固定集合。这一不足正是第三章CADS要解决的核心问题：将计算预算作为数据选择的显式变量。


\section{模型冗余与效率优化}\label{sec:model-redundancy}

\subsection{结构化剪枝}

结构化剪枝（Structured Pruning）通过移除整个滤波器或通道来压缩模型\cite{he2018soft, jiang2022channel}，其产生的稀疏模型可直接在标准硬件上加速推理，无需专用稀疏计算库的支持。

在剪枝准则方面，研究者们提出了多种评估滤波器重要性的方法。基于范数的方法\cite{han2015learning}以滤波器权重的L1或L2范数作为重要性度量，范数较小的滤波器被认为对输出贡献较小，可优先移除。基于梯度敏感度的方法利用Taylor展开近似评估移除某个滤波器对损失函数的影响\cite{lecun1989optimal, hassibi1992second}。此外，还有基于特征图重建误差的方法\cite{he2018soft}，以及基于图卷积网络的通道剪枝方法\cite{jiang2022channel}和自监督掩码学习方法\cite{elkerdawy2022fire}等。

\subsection{非结构化剪枝与稀疏训练}

非结构化剪枝（Unstructured Pruning）在权重级别将单个参数置零\cite{han2015learning, han2016deep}，可达到更高的稀疏率，但产生的不规则稀疏模式需要专用硬件或软件库的支持才能实现实际加速。

经典方法中，基于幅度的剪枝（Magnitude Pruning）\cite{han2015learning}是最直观的方案，直接移除绝对值最小的权重。Frankle和Carlin\cite{frankle2018lottery}提出的彩票假说（Lottery Ticket Hypothesis）揭示了密集网络中存在稀疏子网络，这些子网络从初始化开始独立训练即可达到与密集网络相当的性能。Lee等人\cite{lee2018snip}提出SNIP，利用连接敏感度在训练前一次性确定剪枝掩码。在理论方面，近年来的研究将剪枝与优化理论建立了联系\cite{grosse2016kronecker, zhou2021effective, louizos2017learning, qian2021probabilistic}，为剪枝方法提供了更坚实的理论基础。

稀疏训练（Sparse Training）是一类在训练过程中动态调整稀疏模式的方法。Evci等人\cite{evci2020rigging}提出的RigL方法允许稀疏结构在训练过程中通过"生长"（激活新连接）和"修剪"（移除现有连接）不断演化，而非在训练结束后一次性确定。这种动态调整机制使得稀疏模式能够更好地适应训练过程中模型状态的变化。本文第四章采用RigL作为在线剪枝方法，���是因为其动态特性与交替优化框架天然兼容。

\subsection{小结：数据侧与模型侧"选择"的割裂}

从"冗余消除"的统一视角审视，核心子集选择和模型剪枝本质上是对训练流程中不同维度冗余的"选择"：前者消除数据冗余，识别对训练贡献小的样本；后者消除参数冗余，识别对推理贡献小的权重。值得注意的是，在经典机器学习中，特征筛选（类比于权重剪枝）与样本筛选（类比于核心子集选择）的联合优化已有理论框架\cite{shibagaki2016simultaneous, zhang2017scaling}，经验表明特征筛选可增强样本重要性估计，反之亦然。

然而，在深度学习中，现有工作将这两类选择完全独立进行：先做数据选择再在选出的子集上训练完整模型，或先训练完整模型再做剪枝。这种流水线式处理隐含假设数据冗余与参数冗余相互独立，忽视了二者之间的潜在耦合。这一不足正是第四章SWaST要解决的核心问题：揭示数据冗余与参数冗余的双向耦合关系，提出协同优化方案。


\section{数据选择中的优化方法}\label{sec:optimization-methods}

\subsection{双层优化的基本框架}

双层优化（Bilevel Optimization）是一类具有层次结构的优化问题：外层（上层）目标优化上层变量，内层（下层）目标求解下层变量，且内层的最优解依赖于外层的决策。这种嵌套结构天然适合建模机器学习中的多层次决策问题。

在机器学习中，双层优化的典型应用包括：超参数优化\cite{franceschi2017forward, franceschi2018bilevel, lorraine2020optimizing}，其中外层目标为验证集损失，内层目标为训练集损失，通过联合优化实现模型参数与超参数的协同调优；神经架构搜索，如DARTS\cite{liu2018darts}将架构参数设为外层变量、网络权重设为内层变量，实现高效的架构发现；以及元学习，其中外层优化跨任务的泛化能力，内层完成任务内的快速适应。

\subsection{主要求解策略}

双层优化的核心难点在于计算外层目标关于外层变量的梯度（即超梯度），因为外层目标通过内层最优解间接依赖于外层变量。现有求解策略可分为四类。

\textbf{基于隐式微分的方法}\cite{domke2012generic, grazzi2020iteration, pedregosa2016hyperparameter}利用内层最优性��件（如KKT条件），通过隐式函数定理计算超梯度。这类方法需要计算Hessian向量积，计算代价高且需要内层求解到最优。

\textbf{基于迭代微分的方法}\cite{maclaurin2015gradient, shaban2019truncated}对内层训练过程做截断反向传播（truncated backpropagation），将超梯度表示为训练轨迹上的梯度累积。计算开销与截断步数成正比，在长训练过程中不可行。

\textbf{基于有限差分的方法}\cite{sow2021convergence, yang2023achieving}用参数扰动估计梯度，避免了显式的二阶微分，但引入的方差较大，收敛速度受限。

\textbf{基于惩罚的一阶方法}\cite{chen2023near, kwon2023fully, ye2022bome}是近年来最重要的进展之一。这类方法将内层最优性条件转化为外层目标的惩罚项，消除双层嵌套结构，使得整个问题可以仅用标准一阶梯度求解。这一进展使双层优化算法得以应用于日益大规模的问题\cite{pan2024scalebio}。

\subsection{小结：双层优化在数据选择中的应用与不足}

数据选择问题天然适合双层优化建模：外层优化数据选择策略（使验证集性能最优），内层在选定的数据子集上训练模型\cite{zhou2022probabilistic, killamsetty2021glister, borsos2020coresets}。在这一框架下，数据重加权、核心子集选择以及大规模语言模型的数据配比\cite{pan2024scalebio}等问题都可以统一建模。

然而，现有的双层优化数据选择方法存在一个共性假设：它们要么假设内层训练到收敛（从而可以利用隐式函数定理），要么不显式考虑计算预算约束。当内层因预算限制无法收敛时——这在实际���用中极为常见——标准双层优化的理论基础（如内层最优性条件）不再成立，基于隐式微分的超梯度计算也随之失效。这一不足直接引出第三章的工作：CADS通过惩罚松弛和一维损失近似来处理内层非收敛的情况，将计算预算显式纳入数据选择的优化框架。


\section{分布外检测}\label{sec:ood-detection}

\subsection{问题定义与方法分类}

分布外检测（Out-of-Distribution Detection, OOD Detection）\cite{yang2021generalized_survey}的目标是：给定在分布内（In-Distribution, ID）数据上训练的分类器$f$，设计评分函数$S(x; f)$，使得ID样本的得分系统性地高于（或低于）OOD样本的得分，从而通过阈值$\tau$实现二分判别——$S(x) \geq \tau$则接受输入，否则拒绝。

现有OOD检测方法可从方法论角度分为三大类\cite{yang2021generalized_survey}：（1）基于分类器输出的方法，利用softmax概率\cite{hendrycks2016msp}、logits或能量函数\cite{liu2020energy}等网络输出信息；（2）基于密度估计的方法\cite{zong2018deep_densi, nalisnick2018deep_densi, zisselman2020deep_densi, kirichenko2020normalizing_densi}，在特征空间中显式建模ID数据的概率分布；（3）基于距离的方法\cite{lee2018maha_dist, sun2022knn, ming2022cider_dist}，计算测试���本与ID数据在特征空间中的距离。其中，基于分类器输出的方法通常表现最优\cite{yang2021generalized_survey}。

本文聚焦后处理方法（post-hoc methods），即不修改模型训练过程、不需要OOD数据参与训练，仅在推理时利用已训练模型的输出设计评分函数的方法。这一约束具有重要的实际意义——在真实部署场景中，模型通常不允许因检测需求而重新训练。

\subsection{基于输出层的后处理方法}

最大Softmax概率（Maximum Softmax Probability, MSP）\cite{hendrycks2016msp}是OOD检测最基础的基线方法，它直接以模型输出的最大softmax概率作为ID置信度。MSP简单有效，但由于softmax函数的归一化特性，即使对OOD输入也可能产生较高的置信度，限制了其区分能力。

能量评分（Energy Score）\cite{liu2020energy}用LogSumExp函数替代softmax归一化，将网络输出映射为一个标量能量值。Liu等人\cite{liu2020energy}从理论上证明能量评分与输入的对数似然相关，因此比MSP具有更强的理论支撑和实证表现。

ODIN\cite{liang2018odin_confi}通过两个技术手段增强softmax概率的区分度：温度缩放（temperature scaling）降低softmax输出的尖锐程度，输入扰动（input perturbation）沿梯度方向微调输入以拉大ID与OOD的评分差距。此外，LINe\cite{ahn2023line}通过计算Shapley值减少神经元引入的噪声，GEN\cite{liu2023gen}探索了softmax方法的理论极限。

这些方法的共同特点是均利用网络的最终输出（logits或softmax概率），依赖的信息来源较为单一。

\subsection{基于中间特征的方法}

为了利用比最终输出更丰富的中间表示，一系列方法将检测信号的来源扩展到网络的中间层特征。

\textbf{特征激活操作方法}通过干预中间层的特征激活来改善OOD检测。ReAct\cite{sun2021react}观察到OOD数据倾向于在倒数第二层产生异常大的激活值，通过截断这些异常激活来抑制OOD样本的虚假高分。DICE\cite{sun2022dice}基于单元贡献度对全连接层权重进行稀疏化，提升OOD检测的灵敏度。ASH\cite{djurisic2022extremely}在特征层执行激活整形（activation shaping），通过选择性剪枝策略实现了优于DICE的检测效果。Yu等人\cite{yu2023block}通过选择特征范数区分度最优的网络块来提升检测性能。

\textbf{特征空间距离方法}在特征空间中度量测试样本与ID数据之间的距离。马氏距离方法\cite{lee2018maha_dist}在各层特征空间中计算测试样本到类条件高斯分布中心的马氏距离，综合多层信息进行判断。K近邻方法（KNN）\cite{sun2022knn}直接利用测试样本与ID训练数据在特征空间中的K近邻距离作为OOD评分，无需假设特征服从特定分布。

这些方法利用了比最终输出更丰富的中间表示信息，但一个值得注意的共同点是：大多数方法仍在全局平均池化（Global Average Pooling, GAP）之后操作——无论是特征激活操作还是距离计算，使用的都是池化后的$C$维特征向量。

\subsection{小结：池化前信息的缺失}

全局平均池化将$C \times H \times W$的特征图压缩为$C$维向量，这一操作在降低维度的同时，不可避免地丢失了每个通道内的空间激活分布信息。具体而言，池化后的通道特征值仅反映该通道的平均激活水平，无法区分"少量空间位置产生高激活"与"大量空间位置产生中等激活"这两种截然不同的激活模式。

现有后处理方法——无论是基于输出层的MSP\cite{hendrycks2016msp}、Energy\cite{liu2020energy}，还是基于中间特征的ReAct\cite{sun2021react}、DICE\cite{sun2022dice}、ASH\cite{djurisic2022extremely}、KNN\cite{sun2022knn}——主要在池化后的特征上工作，未利用池化前通道内的激活分布模式。这一信息空白正是第五章NAP的切入点：NAP发现ID样本和OOD样本在池化前的通道内激活模式存在显著差异——ID样本倾向于在少量空间位置产生尖锐的高激活，而OOD样本的激活更为分散。这一"神经激活先验"可作为独立于现有方法的新检测信号，与现有方法实现正交互补。

\chapter{计算预算感知的训练数据选择}

\section{章节引言}

第二章指出，现���训练数据选择方法隐含假设最优数据子集独立于计算预算。本章从一个直觉出发挑战这一假设：在只能训练10轮和可以训练100轮时，"最有价值的数据"应当是相同的吗？

直觉上，答案是否定的。当计算预算有限时，模型仅有时间学习数据中的简单模式\cite{rahaman2018spectral}，此时复杂或含噪声的样本不仅无法被有效利用，反而可能干扰学习过程；当计算预算充裕时，模型有能力从多样化甚至含噪声的数据中提取更丰富的信息，此时仅使用简单数据反而限制了模型的最终能力。这意味着最优数据子集应当依赖于可用的计算预算。

为验证这一直觉，本文设计了一组探索性实验。在由不同噪声水平构成的异构多源数据上，固定数据选择策略（如始终选择最干净的数据源，或始终均匀混合所有源），然后在不同计算预算下评估各策略的表现。图~\ref{fig:cads-exploratory}展示了实验结果。% [VERIFY: 实验设置——CIFAR-10, 5个噪声水平的数据源]

% [图3-1: 不同计算预算下低频与高频数据训练的验证损失对比]
\begin{figure}[htb]
  \centering
  % \includegraphics[width=0.8\textwidth]{figures/cads-exploratory.pdf}
  \fbox{\parbox{0.8\textwidth}{\centering\vspace{2.5cm}
    \textbf{占位图：不同计算预算下低频与高频数据训练的验证损失对比}\\[0.5em]
    对应原论文Fig.1，需手绘
  \vspace{2.5cm}}}
  \caption{不同计算预算下低频与高频数据训练的验证损失对比}
  \label{fig:cads-exploratory}
\end{figure}

实验结果清晰地展示了两个关键现象：第一，没有任何固定数据选择策略能在所有计算预算下一致领先——低预算下偏好干净数据的策略表现最优，而高预算下纳入更多数据源的策略反而更优；% [VERIFY: 具体数字]
第二，不同计算预算下的最优数据配比截然不同。这些发现指向一个此前被忽视的结论：\textbf{最优数据选择策略是计算预算的函数}，预算不应被隐含假设为充裕，而应作为数据选择的显式输入。

基于上述观察，本章提出CADS（Computation-Aware Data Selection）框架，其核心思想是将计算预算$C$作为双层优化中的显式约束，让数据选择策略自动适应给定预算。图~\ref{fig:cads-framework}展示了CADS框架的总体架构。

% [图3-2: CADS框架总览图]
\begin{figure}[htb]
  \centering
  % \includegraphics[width=0.9\textwidth]{figures/cads-framework.pdf}
  \fbox{\parbox{0.9\textwidth}{\centering\vspace{2.5cm}
    \textbf{占位图：CADS框架总览图}\\[0.5em]
    展示双层优化结构、惩罚松弛、CADS-E/S两种变体的关系\\
    需手绘
  \vspace{2.5cm}}}
  \caption{CADS框架总览——双层优化结构、惩罚松弛与两种变体}
  \label{fig:cads-framework}
\end{figure}

实现这一思想面临两个关键技术挑战：（1）预算约束导致内层训练无法收敛，使得标准双层优化方法（如基于隐式微分的方法）失效；（2）数据选择过程本身不应消耗过多计算预算，否则将违背"高效利用数据"的初衷。本章将依次介绍CADS如何应对这两个挑战。


\section{问题定义}

\subsection{计算约束下的数据选择建模}

考虑训练集$\mathcal{D}_{\text{trn}} = \{(x_i, y_i)\}_{i=1}^{N}$和验证集$\mathcal{D}_{\text{val}}$。计算预算$C$以前向-反向传播的总次数度量。数据选择掩码$\mathbf{m} \in \{0,1\}^N$指示哪些样本被选入训练子集。

在计算预算$C$的约束下，数据选择问题可建模为如下双层优化：
\begin{equation}\label{eq:cads-bilevel}
  \min_{\mathbf{s}} \mathbb{E}_{\mathbf{m} \sim p(\mathbf{m}|\mathbf{s})}[\Lval(\btheta_C(\mathbf{m}))]
\end{equation}
其中$\btheta_C(\mathbf{m})$表示在子集$\mathbf{m}$上训练恰好$C$步后得到的模型参数，$p(\mathbf{m}|\mathbf{s})$是以$\mathbf{s}$为参数的采样分布。外层目标是找到使验证损失最小的选择策略$\mathbf{s}$。

这一建模与标准核心子集选择问题\cite{zhou2022probabilistic}的关键区别在于对$\btheta_C(\mathbf{m})$的定义：标准方法假设内层训练到最优$\btheta^*(\mathbf{m})$，而CADS显式将内层训练限制在$C$步。当$C \to \infty$时，CADS退化为标准核心子集选择问题。

\subsection{预算约束带来的根本困难}

标准双层优化假设内层求解到最优$\btheta^*(\mathbf{m})$，在此条件下可利用隐式函数定理\cite{domke2012generic, grazzi2020iteration}计算外层梯度。然而，预算约束下内层仅运行$C$步，$\btheta_C(\mathbf{m}) \neq \btheta^*(\mathbf{m})$，内层最优性条件不成立，隐式微分的理论前提被破坏。

一种替代方案是迭代微分\cite{maclaurin2015gradient, shaban2019truncated}——对$C$步训练过程做截断反向传播。然而，这要求存储整个训练轨迹的中间状态，内存和计算开销均与$C$成正比\cite{chen2016training}，对大预算场景不可行。

因此，CADS需要一种不依赖内层最优性、也不需要对训练过程做梯度回传的新方法来估计外层梯度。


\section{方法：CADS框架}

\subsection{双层优化建模与策略梯度基线}

\textbf{概率化重参数化。}CADS首先将离散掩码$\mathbf{m}$的组合优化问题转化为连续参数$\mathbf{s}$的优化问题。定义参数化采样分布$p(\mathbf{m}|\mathbf{s})$，外层目标变为关于$\mathbf{s}$的期望优化。这一重参数化将搜索空间从$2^N$个候选子集降低为对连续参数$\mathbf{s}$的梯度优化。

\textbf{策略梯度估计。}利用REINFORCE策略梯度\cite{zhou2022probabilistic}，外层梯度可以表示为：
\begin{equation}\label{eq:cads-reinforce}
  \nabla_{\mathbf{s}}\mathbb{E}[\Lval] = \mathbb{E}_{\mathbf{m} \sim p(\mathbf{m}|\mathbf{s})}[\Lval(\btheta_C(\mathbf{m})) \cdot \nabla_{\mathbf{s}} \log p(\mathbf{m}|\mathbf{s})]
\end{equation}
这一估计具有关键优势：它无需对内层训练过程做梯度回传（即无需Hessian计算），外层梯度仅需评估验证损失$\Lval$和采样分布的对数梯度$\nabla_{\mathbf{s}} \log p(\mathbf{m}|\mathbf{s})$。从强化学习的视角看，数据选择$\mathbf{m}$是"动作"，验证损失是"奖励"，$\mathbf{s}$是"策略参数"。

\textbf{Bilevel-CADS的局限。}基于策略梯度的方法（称为Bilevel-CADS）在概念上验证了预算感知数据选择的有效性，但其计算开销仍然很高：每次外层更新需要$K$次独立的完整训练——采样$K$个不同的$\mathbf{m}$，并各自训练$C$步——总计算量为$K \times C$。即使$K$仅取5，在中等预算下也需要5倍完整训练的开销，实际应用中难以接受。Bilevel-CADS作为概念验证，证明了预算感知数据选择的价值，但距离实用化尚有差距。

\subsection{惩罚松弛与单层目标}

为克服Bilevel-CADS的效率瓶颈，CADS提出将双层优化转化为单层优化的惩罚松弛方法\cite{ye2022bome, kwon2023fully}。

\textbf{核心思想。}引入惩罚项，构造如下单层目标：
\begin{equation}\label{eq:cads-penalty}
  \mathcal{L}^{\alpha}_{\text{CADS}} = \mathbb{E}_{p(\mathbf{m}|\mathbf{s})}[\Lval(\btheta) + \alpha(\Ltrn(\btheta, \mathbf{m}) - l(|\mathbf{m}|))^2]
\end{equation}
其中$\alpha$为惩罚系数，$l(|\mathbf{m}|)$为"在大小为$|\mathbf{m}|$的子集上训练$C$步后的最小可达训练损失"。惩罚项的直觉含义是：模型在当前子集上的训练损失应接近该子集大小所对应的"目标损失"——即模型应充分利用所选数据，既不欠拟合（训练损失远高于$l(|\mathbf{m}|)$）也不过拟合（训练损失远低于$l(|\mathbf{m}|)$）。

这一松弛将双层嵌套结构消解为单层联合优化问题：模型参数$\btheta$和选择策略参数$\mathbf{s}$可通过标准SGD交替更新，无需内外层迭代。

\textbf{一维损失近似$l(|\mathbf{m}|)$。}惩罚项中的$l(|\mathbf{m}|)$——即在给定子集大小下$C$步训练后的最小可达损失——如果直接求解，需要对每个候选子集大小做完整训练，计算上不可行。

CADS的一个关键发现是：$l(|\mathbf{m}|)$关于子集大小$|\mathbf{m}|$近似呈对数线性关系（log-linear）。这意味着只需在3至5个不同子集大小上做短训练、记录最终训练损失，即可拟合一条对数线性曲线来近似$l(\cdot)$。% [VERIFY: 拟合细节和精度]
图~\ref{fig:cads-loss-approx}展示了这一近似的效果。

% [图3-3: 一维损失近似的拟合效果]
\begin{figure}[htb]
  \centering
  % \includegraphics[width=0.6\textwidth]{figures/cads-loss-approx.pdf}
  \fbox{\parbox{0.6\textwidth}{\centering\vspace{2cm}
    \textbf{占位图：一维损失近似的拟合效果}\\[0.5em]
    不同子集大小下的实际训练损失与拟合曲线对比\\
    对应原论文Fig.6/附录，需手绘
  \vspace{2cm}}}
  \caption{一维损失近似——不同子集大小下的实际训练损失与对数线性拟合曲线}
  \label{fig:cads-loss-approx}
\end{figure}

这一近似将惩罚项的计算从"每个候选子集独立训练"降低为"一次性拟合 + 查表"，是CADS效率提升的关键所在。

\textbf{效率分析。}CADS的总额外计算开销为$(5/A + 2\gamma)C$，其中$A$为摊销因子（一次选择可复用于多次训练���，$\gamma < 1$为验证集评估比例。% [VERIFY: 具体公式]
作为对比，表~\ref{tab:cads-cost}将CADS与主流数据选择方法的计算开销进行了比较：基于训练动态的方法（如Forgetting\cite{toneva2018empirical}、GraNd/EL2N\cite{paul2021deep}）需要$1.2C$至$2C$的额外开销，而更精细的方法（如Influence Functions、Data Shapley）需要$10C$至$101C$以上的开销。% [VERIFY: Table 6数据]
CADS在大预算场景下可低至约$2\gamma C$的额外开销，且其摊销特性意味着一次选择可复用于多次独立训练。

\subsection{CADS-E：样本级选择}

CADS-E是CADS框架的样本级变体，为每个训练样本独立建模选择概率。

在采样分布设计上，每个样本$i$对应一个独立的Bernoulli参数$s_i \in [0,1]$，表示该样本被选入子集的概率。采样分布为���
\begin{equation}\label{eq:cads-e}
  p(\mathbf{m}|\mathbf{s}) = \prod_{i=1}^{N} s_i^{m_i}(1-s_i)^{1-m_i}
\end{equation}
对数梯度$\partial \log p / \partial s_i = m_i/s_i - (1-m_i)/(1-s_i)$，计算简单且可并行。

CADS-E的优势在于细粒度控制：每个样本拥有独立的选择概率，可识别出特定的噪声样本或特别有价值的样本。其局限在于参数量等于训练集大小$N$，对大规模数据集存在优化效率问题。因此，CADS-E更适用于中小规模数据集或需要样本级精细控制的场景。

\subsection{CADS-S：来源级选择}

CADS-S是CADS框架的来源级变体，将训练数据按来源分组，在来源级别而非样本级别进行选择。

假设训练数据由$n$个来源$\{\mathcal{D}_j\}_{j=1}^{n}$组成，每个来源$j$的采样比例$r_j$服从截断高斯分布：
\begin{equation}\label{eq:cads-s}
  r_j \sim \mathcal{N}(s_j, \sigma_t^2), \quad r_j \in [0,1]
\end{equation}
其中$s_j$为可学习的中心参数，$\sigma_t = 0.99^t \sigma_0$为随训练步数$t$递减的方差。方差退火机制使得选择策略在训练初期保持充分探索，在后期逐渐稳定。

CADS-S的参数量仅为来源数$n$（通常$n \ll N$），适用于多源异构数据的场景，如不同质量的数据集混合、不同域的数据配比。在这类场景中，关键决策是各来源的混合比例，而非逐个样本的取舍。

CADS-S揭示了CADS框架最重要的发现之一：数据选择策略随计算预算系统性地变化。在CIFAR-10异构数据源实验中，低预算时CADS-S几乎只选择最干净的数据源，高噪声源的采样权重趋近于零；随着预算增加，CADS-S逐步纳入中等噪声水平的数据源；极高噪声源（噪声率$>$60\%）在所有预算下均被排除。% [VERIFY: 具体阈值]
这一行为直接体现了数据选择与计算预算的交互：预算决定了模型能"消化"多大难度的数据。


\section{实验}

\subsection{实验设置}

本章设计了四组实验，覆盖不同规模和模态的场景：MNIST\cite{lecun1998mnist}用于小规模概念验证，CIFAR-10\cite{krizhevsky2009cifar}用于异构数据源的系统评估（构造5个不同噪声水平的数据源），DomainNet\cite{domainnet}用于大规模多域场景（6个视觉域），GPT-2\cite{radford2019language}指令微调用于语言模态验证。% [VERIFY: 各实验的详细设置]

基线方法包括：随机选择、均匀混合（Uniform）、最优单源（Best-source）、概率双层核心子集选择（PBCS\cite{zhou2022probabilistic}）以及全数据训练。评价指标为分类准确率（视觉任务）和困惑度（语言任务）。

\subsection{小规模验证}

首先在MNIST数据集上验证Bilevel-CADS基线的概念有效性。表~\ref{tab:cads-mnist}展示了不同数据选择方法在不同初始子集大小下的准确率对比。

% [表3-1: MNIST上不同数据选择方法在不同初始子集大小下的准确率对比]
\begin{table}[htb]
  \centering
  \caption{MNIST上不同数据选择方法在不同初始子集大小下的准确率对比}
  \label{tab:cads-mnist}
  \fbox{\parbox{0.8\textwidth}{\centering\vspace{1.5cm}
    \textbf{占位表：对应原论文Table 1}\\[0.5em]
    % [VERIFY: Table 1]
  \vspace{1.5cm}}}
\end{table}

Bilevel-CADS在不同初始子集大小（200、400、600、800）下均优于随机选择和PBCS，平均准确率达到92.80\%。% [VERIFY: Table 1]
图~\ref{fig:cads-convergence}展示了Bilevel-CADS优化过程中子集大小的收敛轨迹——无论从何种初始大小出发，子集大小均收敛到相近的最终值，表明优化过程对初始化不敏感。

% [图3-4: Bilevel-CADS在不同初始子集大小下的子集大小收敛过程]
\begin{figure}[htb]
  \centering
  % \includegraphics[width=0.6\textwidth]{figures/cads-convergence.pdf}
  \fbox{\parbox{0.6\textwidth}{\centering\vspace{2cm}
    \textbf{占位图：子集大小收敛过程}\\[0.5em]
    对应原论文Fig.2，需手绘
  \vspace{2cm}}}
  \caption{Bilevel-CADS在不同初始子集大小下的子集大小收敛过程}
  \label{fig:cads-convergence}
\end{figure}

这一实验的意义在于证明预算感知的数据选择确实能提升模型性能，为后续高效方法的开发提供了动机基础。

\subsection{异构数据源实验}

\textbf{CIFAR-10实验。}在CIFAR-10上构造由5个不同噪声水平（0\%至90\%）的数据源组成的异构数据集，计算预算从90,000到225,000样本不等。表~\ref{tab:cads-cifar}展示了各方法在不同预算下的准确率。

% [表3-2: CIFAR-10异构数据源实验]
\begin{table}[htb]
  \centering
  \caption{CIFAR-10异构数据源实验——不同计算预算下各方法的准确率对比}
  \label{tab:cads-cifar}
  \fbox{\parbox{0.8\textwidth}{\centering\vspace{1.5cm}
    \textbf{占位表：对应原论文Table 2}\\[0.5em]
    % [VERIFY: Table 2]
  \vspace{1.5cm}}}
\end{table}

CADS-S平均准确率63.08\%，显著优于Best-source基线的61.65\%和全数据训练。% [VERIFY: Table 2]
更重要的是，CADS-S在所有计算预算下一致领先，而其他方法的相对优劣随预算变化。图~\ref{fig:cads-cifar-ratio}展示了CADS-S在四种不同预算下的采样比例演化过程，清晰地展示了选择策略如何随预算自适应调整。

% [图3-5: CIFAR-10上不同计算预算下CADS-S的采样比例演化过程]
\begin{figure}[htb]
  \centering
  % \includegraphics[width=\textwidth]{figures/cads-cifar-ratio.pdf}
  \fbox{\parbox{\textwidth}{\centering\vspace{2cm}
    \textbf{占位图：CIFAR-10采样比例演化过程}\\[0.5em]
    四个子图分别对应四种预算，对应原论文Fig.3，需手绘
  \vspace{2cm}}}
  \caption{CIFAR-10上不同计算预算下CADS-S的采样比例演化过程}
  \label{fig:cads-cifar-ratio}
\end{figure}

\textbf{DomainNet实验。}在包含6个视觉域的大规模DomainNet数据集上，CADS-S进一步验证了其可扩展性。表~\ref{tab:cads-domainnet}展示了结果。

% [表3-3: DomainNet上不同初始采样比例下各方法的准确率对比]
\begin{table}[htb]
  \centering
  \caption{DomainNet上不同初始采样比例下各方法的准确率对比}
  \label{tab:cads-domainnet}
  \fbox{\parbox{0.8\textwidth}{\centering\vspace{1.5cm}
    \textbf{占位表：对应原论文Table 3}\\[0.5em]
    % [VERIFY: Table 3]
  \vspace{1.5cm}}}
\end{table}

CADS-S平均准确率42.27\%，远优于均匀采样基线的36.22\%，最高提升达14.42\%。% [VERIFY: Table 3]
图~\ref{fig:cads-domainnet-ratio}展示了DomainNet上的采样比例演化。

% [图3-6: DomainNet上CADS-S的采样比例演化过程]
\begin{figure}[htb]
  \centering
  % \includegraphics[width=0.8\textwidth]{figures/cads-domainnet-ratio.pdf}
  \fbox{\parbox{0.8\textwidth}{\centering\vspace{2cm}
    \textbf{占位图：DomainNet采样比例演化过程}\\[0.5em]
    对应原论文Fig.4，需手绘
  \vspace{2cm}}}
  \caption{DomainNet上CADS-S的采样比例演化过程}
  \label{fig:cads-domainnet-ratio}
\end{figure}

这两组实验共同验证了核心发现：CADS-S能自动学习到适配当前计算预算的数据配比策略，在所有预算设定下一致优于固定策略基线。

\subsection{指令微调实验}

为验证CADS在语言模态上的通用性，本文在GPT-2\cite{radford2019language}指令微调任务上进行了实验。

\textbf{双源实验。}使用Alpaca-GPT4\cite{peng2023instruction}（1K条高质量指令）和Alpaca\cite{alpaca}（9K条标准质量指令）作为两个数据源，计算预算从$1\times10^4$到$5\times10^4$。表~\ref{tab:cads-alpaca}展示了各方法在不同预算下的困惑度对比。

% [表3-4: GPT-2指令微调——不同计算预算下各方法的困惑度对比]
\begin{table}[htb]
  \centering
  \caption{GPT-2指令微调——不同计算预算下各方法的困惑度对比}
  \label{tab:cads-alpaca}
  \fbox{\parbox{0.8\textwidth}{\centering\vspace{1.5cm}
    \textbf{占位表：对应原论文Table 4}\\[0.5em]
    % [VERIFY: Table 4]
  \vspace{1.5cm}}}
\end{table}

CADS-S在困惑度指标上一致优于基线方法。% [VERIFY: Table 4]
图~\ref{fig:cads-alpaca-ratio}展示了两个数据源的采样比例随预算的动态变化，再次体现了预算对最优数据配比的影响。

% [图3-7: Alpaca数据集上不同计算预算下的采样比例动态]
\begin{figure}[htb]
  \centering
  % \includegraphics[width=0.7\textwidth]{figures/cads-alpaca-ratio.pdf}
  \fbox{\parbox{0.7\textwidth}{\centering\vspace{2cm}
    \textbf{占位图：Alpaca采样比例动态}\\[0.5em]
    对应原论文Fig.5，需手绘
  \vspace{2cm}}}
  \caption{Alpaca数据集上不同计算预算下的采样比例动态}
  \label{fig:cads-alpaca-ratio}
\end{figure}

\textbf{大规模数据混合实验。}进一步在13个指令微调数据集\cite{alpaca, peng2023instruction, lian2023slimorca, openorcagpt4, gpteacher, khashabi2020unifiedqa, liu2019multi, wei2021finetuned, bukharin2023data, raffel2019exploring}上验证CADS-S的可扩展性。表~\ref{tab:cads-13source}展示了结果。

% [表3-5: 13个指令微调数据集的大规模数据混合实验结果]
\begin{table}[htb]
  \centering
  \caption{13个指令微调数据集的大规模数据混合实验结果}
  \label{tab:cads-13source}
  \fbox{\parbox{0.8\textwidth}{\centering\vspace{1.5cm}
    \textbf{占位表：对应原论文Table 5}\\[0.5em]
    % [VERIFY: Table 5]
  \vspace{1.5cm}}}
\end{table}

在13个数据源的混合场��中，CADS-S仍然能够有效地优化数据配比，证明了其在大规模多源数据选择中的实用价值。% [VERIFY: Table 5]

\subsection{计算开销分析}

表~\ref{tab:cads-cost}将CADS与主流数据选择方法的计算开销进行了系统对比。

% [表3-6: 各数据选择方法的计算开销对比分析]
\begin{table}[htb]
  \centering
  \caption{各数据选择方法的计算开销对比分析}
  \label{tab:cads-cost}
  \fbox{\parbox{0.8\textwidth}{\centering\vspace{1.5cm}
    \textbf{占位表：对应原论文Table 6}\\[0.5em]
    % [VERIFY: Table 6]
  \vspace{1.5cm}}}
\end{table}

对比结果表明，CADS的计算开销随预算规模增大而摊薄。% [VERIFY: Table 6]
在大预算场景下，CADS的额外开销远低于基于训练动态的方法（Forgetting、GraNd、EL2N\cite{toneva2018empirical, paul2021deep}），更远低于基于数据价值估计的方法（Influence Functions\cite{pruthi2020estimating}、Data Shapley）。这一效率优势使CADS成为大规模数据选择场景的实用方案。


\section{讨论}

\subsection{损失预测器的角色与局限性}

CADS的惩罚松弛方法依赖一维损失近似$l(|\mathbf{m}|)$的准确性。如果该近似偏差过大，惩罚项将引导模型进入错误的训练状态，导致选择策略失效。

对数线性近似在以下情况下可能不够准确：数据分布极端不均匀、计算预算范围跨越多个数量级、数据来源间质量差异极大。在这些极端场景下，可能需要更灵活的拟合方法，如分段线性拟合或使用Scaling Law中常用的四参数函数族\cite{kaplan2020scaling, hoffmann2022training}。

\subsection{与大模型数据混合领域的联系}

CADS-S的来源级选择与LLM预训练中的数据混合（Data Mixture）问题在形式上高度一致：二者都需要确定不同数据来源的最优混合比例\cite{albalak2024survey}。LLM领域的损失预测器和Scaling Law方法（如Chinchilla\cite{hoffmann2022training}）通过预测不同配比下的损失来指导数据混合，与CADS中的一维损失近似$l(|\mathbf{m}|)$扮演类似角色。

二者的关键异同在于：CADS显式引入计算预算作为变量，使得选择策略随预算自适应变化；而大多数LLM数据混合工作假设固定的预算（固定token数）。CADS的预算感知思路可能对LLM领域的数据配比研究提供新的视角。


\section{章节小结}

本章提出了CADS框架，首次将计算预算纳入数据选择的优化框架，揭示了一个被广泛忽视的交互：\textbf{最优数据子集是计算预算的函数}——低预算偏好简单干净的数据，高预算受益于多样化的数据。在方法层面，CADS通过概率化重参数化的策略梯度估计绕过了内层非收敛的困难，进一步通过惩罚松弛和一维损失近似将昂贵的双层优化转化为高效的单层求解。在实证层面，CADS跨视觉分类和语言模型微调任务验证了预算感知数据选择的一致优势，最高提升14.42\%。% [VERIFY]

CADS在优化数据选择时假设模型结构固定——模型作为"黑盒"接收所选数据并返回验证损失。然而，模型本身也存在冗余参数\cite{han2015learning, frankle2018lottery}，而模型容量（参数量）与数据量之间存在已知的Scaling关系\cite{kaplan2020scaling, hoffmann2022training}。一个自然的追问是：当模型容量也成为可优化变量时——例如通过剪枝压缩模型——数据选择与参数冗余之间是否存在类似的交互？第四章将回答这一问题。

\include{ch4-swast}
\include{ch5-nap}
\chapter{总结与展望}

\section{工作总结}

本文围绕数据选择这一核心方法论，针对深度学习全流程中的三个关键子问题，分别提出了一个框架或方法。

\textbf{CADS：计算预算感知的数据选择（第三章）。}针对训练阶段数据选择与计算资源约束的交互，本文提出CADS框架。CADS首次将计算预算纳入数据选择的双层优化框架，通过概率化重参数化的策略梯度绕过内层非收敛的困难，并通过惩罚松弛与一维损失近似实现高效求解。CADS的核心发现是：\textbf{最优数据子集是计算预算的函数}——低预算偏好简单干净的数据，高预算受益于多样化的数据。在视觉分类和语言模型微调任务上，CADS在所有计算预算下一致领先于基线方法。

\textbf{SWaST：数据选择与权重剪枝的协同优化（第四章）。}针对训练阶段数据选择与模型容量约束的交互，本文提出SWaST框架。SWaST首次透明地分析了数据冗余与参数冗余的双向耦合关系，揭示了独立优化两类冗余的局限性。通过交替执行权重剪枝与核心子集选择，SWaST建立了协同效应：精选数据帮助剪枝（减少过拟合，使权重重要性估计更准确），剪枝后的紧凑模型帮助数据选择（降低特征空间有效维度，使核心子集选择更精准）。SWaST同时识别并解决了"双重损失"这一独特挑战，通过状态保护机制防止不可逆退化。

\textbf{NAP：基于神经激活先验的分布外检测（第五章）。}针对推理阶段数据选择与分布偏移的交互，本文提出NAP方法。NAP发现了一个被现有方法忽视的OOD检测信号：训练使模型在ID数据上形成尖锐的通道内激活先验——少量空间位置产生高激活，其余位置保持低激活。这一先验在OOD样本上不成立（激活更分散），因此可直接用于推理时的OOD判断。NAP设计的即插即用评分函数与现有方法正交互补，在CIFAR和ImageNet基准上与多种基线组合后均实现显著提升。


\section{统一视角}

\subsection{阶段覆盖的完整性}

三项工作覆盖了数据选择在深度学习流程中的三个关键环节：训练阶段的子集选择（CADS）、训练阶段的数据-模型联合选择（SWaST）、推理阶段的数据有效性判断（NAP）。这一覆盖从"选什么数据训练"出发，经过"数据与模型如何协同选择"，到达"如何判断输入是否处于有效域内"，形成了从训练到推理的完整链条。

\subsection{交互视角的统一性}

表~\ref{tab:unified-comparison}从交互维度统一对比三项工作。

% [表6-1: 三项工作的统一对比总结]
\begin{table}[htb]
  \centering
  \caption{三项工作的统一对比——研究问题、阶段定位、交互变量、核心方法与关键发现}
  \label{tab:unified-comparison}
  \fbox{\parbox{0.9\textwidth}{\centering\vspace{2cm}
    \textbf{占位表：三项工作的统一对比总结}\\[0.5em]
    列：工作名称、研究问题、阶段、交互变量、核心方法、关键发现
  \vspace{2cm}}}
\end{table}

三项工作揭示了一个共性的方法论启示：数据选择的最优策略不是孤立的，而是取决于其所处的约束环境。计算资源约束改变了最优训练数据的组成（CADS），模型容量约束与数据冗余相互耦合、需要协同而非独立处理（SWaST），推理时分布偏移要求对数据有效性做在线判断、连接了训练时的选择与推理时的检测（NAP）。

这一共性可以提炼为：\textbf{未来的数据选择研究应系统考虑选择策略与其交互因素的依赖关系，而非将数据选择作为孤立的优化问题。}


\section{局限性}

\subsection{CADS的局限}

CADS的一维损失近似$l(|\mathbf{m}|)$采用的对数线性假设在极端数据分布下可能不够准确，需要更灵活的拟合方法\cite{kaplan2020scaling}。当前实验规模（最大DomainNet约0.5M样本）尚未在十亿级预训练数据上验证，CADS在更大规模上的表现有待探索。此外，CADS在每个计算预算下独立优化选择策略，未利用不同预算间可能存在的共享信息。% [VERIFY: 原论文limitations]

\subsection{SWaST的局限}

状态保护机制引入了额外的前向传播和KL散度计算开销，在极大规模数据集上可能成为瓶颈。当前SWaST仅与RigL\cite{evci2020rigging}剪枝算法验证了兼容性，与其他剪枝方法（如结构化剪枝\cite{he2018soft}）和其他压缩技术（量化\cite{han2016deep, courbariaux2016binarized}、知识蒸馏）的兼容性未经探索。交替频率$\mathcal{R}$和保护强度$\lambda$的调优目前依赖经验搜索，缺乏理论指导。% [VERIFY: 原论文limitations]

\subsection{NAP的局限}

NAP依赖全局池化层（CNN）或注意力机制（Transformer）的存在，对不含此类结构的架构不直接适用。在近分布OOD（near-OOD）场景下——即OOD数据与ID数据在语义上高度相似时——通道内激活模式的差异也会缩小，检测区分度有限。组合权重$w$需要通过伪OOD数据调优，不同ID/OOD对的最优$w$可能不同。% [VERIFY: 原论文limitations]


\section{未来研究方向}

\subsection{跨阶段的数据选择联合建模}

本文三项工作分别处理不同阶段的数据选择，但它们之间存在潜在联系：训练阶段的数据选择决定了模型的有效域，这直接影响推理时OOD检测的难度。一个有意义的未来方向是：能否在训练阶段做数据选择时，直接优化推理时的OOD检测能力？例如，选择能使模型学到更清晰决策边界的训练子集，从而在推理时更容易区分ID与OOD输入。更广义地，端到端考虑从数据选择到模型训练到推理部署的全链条优化，是一个值得探索的方向。

\subsection{大模型时代的数据选择新挑战}

LLM预训练面临万亿级数据的选择与配比问题\cite{albalak2024survey}。CADS的预算感知思路和来源级选择（CADS-S）可为此提供新视角：在不同训练阶段（预训练初期 vs. 后期），最优数据配比是否也应随之变化？LLM微调中的数据选择\cite{zhou2023lima}还需要在数据质量、多样性和安全性（alignment）之间寻求平衡，计算预算约束更为严格。此外，多模态大模型中跨模态数据的配比和互补性选择，也是数据选择领域面临的新挑战。

\subsection{更多交互维度的探索}

本文探索了数据选择与三个交互变量（计算资源、模型容量、推理时分布偏移）的关系。类似的交互可能存在于更多维度：数据选择与模型架构搜索（NAS\cite{liu2018darts}）之间是否存在耦合？数据选择与其他压缩技术（量化\cite{han2016deep}、知识蒸馏）之间是否存在类似SWaST发现的协同关系？在持续学习场景中，数据分布随时间变化\cite{borsos2020coresets}，数据选择策略需要如何适应？在联邦学习场景中，数据分布于多个参与方，又该如何协调选择？这些问题构成了数据选择领域的广阔未来图景。


% ======== 后置部分 ========
\backmatter

\printbibliography

\begin{acknowledgements}
  % TODO: 致谢
  致谢内容。
\end{acknowledgements}

\end{document}
